% ===================================================================
%  TABLEAU N° 15 — Le marché. --- Les aliments.
%
%  Traduction, faite en 2026, en francais standard du Quebec, sur
%  l'ido de texto/io. Regles de la colonne : voir 10-tableau-01.tex.
%
%  LES NEUF CHOIX PROPRES A CE TABLEAU :
%    (24) vente aux enchères ..... vente à l'encan
%    (36) liqueurs ............... spiritueux
%    (45) garçon épicier ......... commis
%    (97) écaillère .............. marchande d'huîtres
%    (86) blanchisseuse-repasseuse . repasseuse
%    (26) ménagères .............. maîtresses de maison
%    (8)  crémière ............... marchande de beurre et de fromage
%    (100) agent de police ....... policier
%         torchons ............... linges à vaisselle
%
%  « ENCAN » EST LE MOT DU QUEBEC, ET C'EST ICI QU'IL SERT. L'ido dit
%  « auciono-vendo », la vente au plus offrant, et c'est l'heure ou la
%  halle s'anime : au Quebec cette vente-la s'appelle un encan, le mot
%  est courant, et « vente aux encheres » y sonne comme un terme de
%  droit. Deux mots de la meme famille — encanteur, encanter — vivent
%  avec lui.
%
%  « LIQUEURS » RETOMBE POUR LA TROISIEME FOIS. Apres les bouteilles
%  du tableau 4 et le liquoriste du tableau 14, l'epicerie du tableau
%  15 vend des « liquori » que l'ido enumere aussitot : rhum, cognac,
%  kirsch, anisette, curacao. Ce sont des spiritueux, et au Quebec la
%  liqueur est une boisson gazeuse.
%
%  « LINGE A VAISSELLE » CONTRE « TORCHON ». Le fac-simile ecrit
%  torchon pour ce que l'ido appelle « vishilo », l'essuie-tout de la
%  cuisine. Au Quebec le torchon est le chiffon sale avec lequel on
%  lave le plancher, et le linge propre qui essuie les assiettes est
%  un linge a vaisselle. La marchande de toile en vend des piles :
%  ce sont donc des linges a vaisselle.
% ===================================================================

%%K t15-apar-1 apar
\VUsaut{4.0mm}
\VUtitre{15.0pt}{60}{TABLEAU N\textsuperscript{o} 15}
\VUsaut{3.0mm}

%%K t15-tit-1 sub
\VUpk{11.4pt}{40}{Le marché. --- Les aliments.}
\VUsaut{3.0mm}

%%K t15-01-1 p
1. --- La plupart des commerçants qui nous fournissent les denrées nécessaires à
notre alimentation se rassemblent sous la \VUgras{halle} \textsuperscript{(1)} du
\VUgras{marché couvert} ou dans les rues voisines.

%%K t15-02-1 p
2. --- Le \VUgras{charcutier} \textsuperscript{(2)}, qui vend du \VUgras{jambon}
\textsuperscript{(3)}, du \VUgras{boudin} \textsuperscript{(4)}, des
\VUgras{saucisses} \textsuperscript{(5)}, des \VUgras{andouillettes}
\textsuperscript{(6)} et du \VUgras{lard} \textsuperscript{(7)}, se tient près de
la \VUgras{marchande de beurre et de fromage} \textsuperscript{(8)}, dont
l'étalage est garni de \VUgras{fromages de Hollande} \textsuperscript{(9)} à
croûte rouge, de \VUgras{camemberts} \textsuperscript{(10)}, de \VUgras{meules de
gruyère} \textsuperscript{(11)} et de \VUgras{mottes de beurre}
\textsuperscript{(12)} frais.

%%K t15-03-1 p
3. --- Devant elle, une \VUgras{marchande de légumes} \textsuperscript{(13)}
propose à ses clientes de belles \VUgras{tomates} \textsuperscript{(14)}, des
\VUgras{choux-fleurs} \textsuperscript{(15)}, de la \VUgras{salade}
\textsuperscript{(16)}, des \VUgras{choux pommés} \textsuperscript{(17)}, des
\VUgras{citrouilles} \textsuperscript{(19)}, des \VUgras{melons}
\textsuperscript{(18)} et même des \VUgras{champignons} \textsuperscript{(20)}.
Mais un \VUgras{livreur de bière}, qui a arrêté sa \VUgras{voiture de livraison}
\textsuperscript{(21)} chargée de \VUgras{tonneaux de bière}
\textsuperscript{(22)} juste devant son étalage, empêche ses clientes
d'approcher. Une jardinière lui apporte un panier d'\VUgras{artichauts}
\textsuperscript{(23)}.

%%K t15-tit-2 sub
\VUsaut{4.0mm}
\VUpk{10.2pt}{40}{Vendre et acheter.}
\VUsaut{3.0mm}

%%K t15-04-1 p
4. --- Au marché, l'animation est grande, parce que l'heure de la \VUgras{vente
à l'encan} \textsuperscript{(24)} est arrivée. Les \VUgras{maîtresses de maison}
\textsuperscript{(26)}, qui viennent y faire leurs achats, s'arrêtent en passant
devant la boutique de la \VUgras{marchande de tripes} \textsuperscript{(25)}
pour acheter des \VUgras{tripes} ou une \VUgras{tête de veau}.

%%K t15-05-1 p
5. --- Un gros \VUgras{cuisinier} \textsuperscript{(27)} marchande un très beau
\VUgras{thon} chez la \VUgras{poissonnière} \textsuperscript{(28)}, qui augmente
ses prix à sa guise. Elle a reçu une grande quantité d'\VUgras{anguilles}, de
\VUgras{soles}, de \VUgras{truites} et de \VUgras{carpes}. Sa \VUgras{morue} et
ses \VUgras{moules} sont très fraîches.

%%K t15-tit-3 sub
\VUsaut{4.0mm}
\VUpk{10.2pt}{40}{Bon marché et cher.}
\VUsaut{3.0mm}

%%K t15-06-1 p
6. --- La \VUgras{marchande de volailles}\textsuperscript{(29)},
installée près d'elle, est très consciencieuse. Quand elle vend une
\VUgras{dinde}, une \VUgras{oie}, un \VUgras{canard} ou un simple
\VUgras{poulet}, elle ne demande que le juste prix.

%%K t15-07-1 p
7. --- Au coin de la place, au premier étage, un \VUgras{marchand de toile}
\textsuperscript{(30)} s'est installé depuis peu; les clientes y sont toujours
sûres de trouver un très beau choix de \VUgras{nappes}, de \VUgras{serviettes},
de \VUgras{tabliers} et de \VUgras{linges à vaisselle}. Au même étage, un
\VUgras{coutelier} \textsuperscript{(31)} a installé son atelier. Il aiguise les
\VUgras{couteaux} des marchandes de la halle et fabrique pour les jardiniers des
\VUgras{serpes} dans l'\VUgras{acier} le plus dur.

%%K t15-tit-4 sub
\VUsaut{4.0mm}
\VUpk{10.2pt}{40}{L'épicerie.}
\VUsaut{3.0mm}

%%K t15-08-1 p
8. --- Au-dessous de son atelier, un grand magasin d'alimentation a ouvert
depuis peu. Ce n'est déjà plus la simple et modeste \VUgras{épicerie}
\textsuperscript{(32)} d'autrefois, qui ne vendait que les denrées courantes :
\VUgras{thé} \textsuperscript{(33)}, \VUgras{chocolat} \textsuperscript{(34)},
\VUgras{pain de sucre} \textsuperscript{(37)} et \VUgras{savon}
\textsuperscript{(38)}. On y vend de tout : des \VUgras{biscuits secs}, des
\VUgras{conserves} \textsuperscript{(35)}, des \VUgras{spiritueux}
\textsuperscript{(36)}, du \VUgras{rhum}, du \VUgras{cognac}, du
\VUgras{kirsch}, de l'\VUgras{anisette} et du \VUgras{curaçao}. On y trouve du
gibier : \VUgras{lièvre} \textsuperscript{(39)}, \VUgras{grives}
\textsuperscript{(40)}, \VUgras{perdrix} \textsuperscript{(41)},
\VUgras{cailles} \textsuperscript{(42)}, \VUgras{canard sauvage}
\textsuperscript{(43)} ou \VUgras{faisan} \textsuperscript{(44)}, aussi
facilement que des fruits exotiques comme les \VUgras{bananes}
\textsuperscript{(46)} et les \VUgras{ananas}.

%%K t15-09-1 p
9. --- Un \VUgras{commis} \textsuperscript{(45)} très serviable est prêt à
donner ce qu'on désire : de la \VUgras{confiture} \textsuperscript{(47)} de
groseilles ou de coings, de la \VUgras{graisse} \textsuperscript{(48)}, des
\VUgras{cornichons} \textsuperscript{(49)} ou des \VUgras{sardines}
\textsuperscript{(50)} salées.

%%K t15-09-2 p
C'est très commode pour les \VUgras{clientes} \textsuperscript{(51)}, qui
trouvent, à côté des denrées les plus courantes, les primeurs les plus rares et
les fruits les plus savoureux : \VUgras{poires} \textsuperscript{(52)}
juteuses, \VUgras{prunes} \textsuperscript{(53)}, \VUgras{raisins}
\textsuperscript{(54)}, \VUgras{pêches} \textsuperscript{(55)},
\VUgras{amandes} \textsuperscript{(56)} et \VUgras{noix}
\textsuperscript{(57)}.

%%K t15-tit-5 sub
\VUsaut{4.0mm}
\VUpk{10.2pt}{40}{La clientèle.}
\VUsaut{3.0mm}

%%K t15-10-1 p
10. --- Le \VUgras{riz} \textsuperscript{(58)} et les \VUgras{lentilles}
\textsuperscript{(59)} sont à côté de la caisse des \VUgras{harengs}
\textsuperscript{(60)} fumés; les \VUgras{pommes de terre} \textsuperscript{(61)}
et les \VUgras{haricots} \textsuperscript{(63)} voisinent avec les
\VUgras{pommes} \textsuperscript{(62)}, les \VUgras{figues}
\textsuperscript{(64)}, les \VUgras{pruneaux} \textsuperscript{(65)} et les
\VUgras{noisettes} \textsuperscript{(66)}. On trouve dans ce magasin des
\VUgras{œufs} \textsuperscript{(67)} frais et des \VUgras{olives}
\textsuperscript{(68)}; aussi les \VUgras{paniers} \textsuperscript{(70)} et les
\VUgras{filets à provisions} \textsuperscript{(73)} se remplissent-ils très vite.

%%K t15-11-1 p
11. --- Le \VUgras{café} \textsuperscript{(72)} de cette excellente maison s'est
acquis, chez les \VUgras{domestiques} \textsuperscript{(69)} et les
\VUgras{cuisinières}, une juste réputation, parce que le vieil
\VUgras{épicier} \textsuperscript{(71)} le torréfie lui-même avec le plus grand
soin.

%%K t15-tit-6 sub
\VUsaut{4.0mm}
\VUpk{10.2pt}{40}{Le spectacle de la rue.}
\VUsaut{3.0mm}

%%K t15-12-1 p
12. --- Tout le monde ne va pas au marché pour s'approvisionner. Dans la
circulation pittoresque de la ruelle qui mène à la halle, on voit les personnes
les plus diverses. À côté d'un aimable \VUgras{garçon boucher}
\textsuperscript{(75)}, qui pousse devant lui une \VUgras{brouette}
\textsuperscript{(74)}, voici deux soldats en permission, tout fiers de montrer
leur uniforme brillant : le \VUgras{hussard} \textsuperscript{(76)} avec son
\VUgras{dolman} bleu et son \VUgras{shako} à plumet, et le \VUgras{caporal de
zouaves}, avec sa culotte bouffante et la tête couverte d'un \VUgras{fez}.

%%K t15-13-1 p
13. --- Le \VUgras{boulanger} \textsuperscript{(80)}, qui se tient sur le seuil
de la \VUgras{boulangerie} \textsuperscript{(78)}, vient de pétrir la pâte de son
\VUgras{pain} \textsuperscript{(79)} et de retirer du four les
\VUgras{croissants} \textsuperscript{(81)}, les \VUgras{petits pains au lait}
\textsuperscript{(82)} et les \VUgras{pains ronds} \textsuperscript{(83)} qui
sont dans la vitrine.

%%K t15-14-1 p
14. --- Au-dessus de la boulangerie habite une habile \VUgras{lingère}
\textsuperscript{(84)} qui emploie plusieurs \VUgras{brodeuses}
\textsuperscript{(85)}. À côté d'elle habite une \VUgras{repasseuse}
\textsuperscript{(86)}, qui empèse le \VUgras{linge} \textsuperscript{(88)} après
l'avoir lavé et fait sécher, et qui le repasse avec son \VUgras{fer à repasser}
\textsuperscript{(87)}.

%%K t15-tit-7 sub
\VUsaut{4.0mm}
\VUpk{10.2pt}{40}{Viande, poissons et légumes.}
\VUsaut{3.0mm}

%%K t15-15-1 p
15. --- À la \VUgras{boucherie} \textsuperscript{(89)}, située au
rez-de-chaussée, on vend toutes sortes de \VUgras{viandes} : du
\VUgras{mouton} \textsuperscript{(90)}, du \VUgras{bœuf} \textsuperscript{(94)}
ou du \VUgras{veau} \textsuperscript{(92)}, en \VUgras{gigots}, en
\VUgras{biftecks}, en \VUgras{rôtis} et en \VUgras{côtelettes}. Le
\VUgras{garçon boucher} \textsuperscript{(93)} découpe toutes ces viandes et met
à part les \VUgras{os à moelle} \textsuperscript{(96)}. À la porte de la
boucherie se tient une \VUgras{marchande d'huîtres} \textsuperscript{(97)} qui
vend des \VUgras{huîtres} \textsuperscript{(98)}. Elle les ouvre si on le
souhaite.

%%K t15-16-1 p
16. --- Tous ces commerçants paient une patente ou une taxe de place; aussi le
\VUgras{policier} \textsuperscript{(100)} poursuit-il sans pitié les pauvres
\VUgras{marchandes de légumes ambulantes} \textsuperscript{(99)}, qui leur font
concurrence en colportant sur leurs voiturettes des \VUgras{carottes}, des
\VUgras{radis}, des \VUgras{navets}, des \VUgras{poireaux} ou des \VUgras{bottes
d'asperges}, des \VUgras{petits pois} frais, des \VUgras{épinards}, de
l'\VUgras{oseille}, du \VUgras{cresson} et du \VUgras{persil}.

%%K t15-tit-8 sub
\VUsaut{4.0mm}
\VUpk{10.2pt}{40}{Attention au policier!}
\VUsaut{3.0mm}

%%K t15-17-1 p
17. --- Le garçon qui vend de l'\VUgras{ail} \textsuperscript{(101)} et des
\VUgras{oignons} sait très bien que lui non plus n'a pas le droit de vendre sa
marchandise près du marché. Il se sauve en courant, au risque de renverser le
panier de \VUgras{crabes}, d'\VUgras{écrevisses} et de \VUgras{crevettes} du
\VUgras{marchand} \textsuperscript{(102)} de \VUgras{langoustes} et de
\VUgras{homards}.

%%K t15-18-1 p
18. --- Tout le monde s'agite et fait grand bruit; mais cela n'empêche pas
d'entendre clairement la voix du \VUgras{vitrier} \textsuperscript{(103)}
ambulant, ni le cri du \VUgras{charbonnier} \textsuperscript{(104)} qui vient de
quitter quelques instants sa charrette pour livrer un \VUgras{sac de charbon}
\textsuperscript{(105)}.
\VUsaut{6.0mm}
\VUfilet{22mm}
